% !TeX root = ../Project.tex
All-atom protein models compress atomic structure into latent representations, often using reconstruction or generative objectives that require these latents to preserve precise coordinate information. Yet, residue properties depend on both local atomic environments and longer-range 3D interactions. Joint-Embedding Predictive Architectures (JEPAs) offer an alternative: predicting representations in latent space rather than reconstructing the raw input. In vision, this shifts learning away from pixel-level reconstruction toward higher-level features that transfer across tasks. We ask whether the same approach can help protein representations capture both residue-level information and broader structural context useful for protein structure and function tasks. Using the La-Proteina encoder, we train JEPA, JEPA+SIGReg and LeJEPA on masked protein structures and evaluate them with StructTokenBench. We study how masking, positional information, rotations and cropped views affect the learned representations. JEPA with 70\% random masking scores higher than the C$\alpha$-conditioned La-Proteina baseline on 16 of 20 functional-site and flexibility scores. In JEPA+SIGReg, I-JEPA-style absolute positional embeddings create a strong residue-index shortcut. Relative positioning suppresses this direct shortcut and scores higher on 18 of 20 functional-site and flexibility outputs. LeJEPA reveals a different shortcut: its representations organise residues by chain length, motivating us to replace full-protein views with fixed-length crops. Evaluating intermediate encoder representations shows that uncropped LeJEPA's block 6 representation outperforms the final representations of all 18 models in our main comparison on 22 of 24 downstream scores. These results show the potential of latent prediction for all-atom protein representations while highlighting the importance of protein-specific representation design.
